\documentclass[10pt,a4paper]{amsart}
\usepackage{amsmath,amssymb,amsthm,graphicx}
% Retained minimal preamble for the Stacks Project chapter sites-cohomology.tex (Cohomology on Sites).
% Verbatim excerpt of the chapter preamble: Xy-pic setup (native preamble lines 14-18) and the native
% theorem declarations with the mathematical macros (native preamble lines 154-217). Copied mechanically;
% nothing added, removed or redefined. The wrapper class and the amsmath/amssymb/amsthm/graphicx base
% packages are supplied by the assembler, so the class/externaldocument part of the native preamble is omitted.
% For commutative diagrams we use Xy-pic
\usepackage[all]{xy}

% We use 2cell for 2-commutative diagrams.
\xyoption{2cell}
% Theorem environments.
%
\theoremstyle{plain}
\newtheorem{theorem}[subsection]{Theorem}
\newtheorem{proposition}[subsection]{Proposition}
\newtheorem{lemma}[subsection]{Lemma}

\theoremstyle{definition}
\newtheorem{definition}[subsection]{Definition}
\newtheorem{example}[subsection]{Example}
\newtheorem{exercise}[subsection]{Exercise}
\newtheorem{situation}[subsection]{Situation}

\theoremstyle{remark}
\newtheorem{remark}[subsection]{Remark}
\newtheorem{remarks}[subsection]{Remarks}

\numberwithin{equation}{subsection}

% Macros
%
\def\lim{\mathop{\mathrm{lim}}\nolimits}
\def\colim{\mathop{\mathrm{colim}}\nolimits}
\def\Spec{\mathop{\mathrm{Spec}}}
\def\Hom{\mathop{\mathrm{Hom}}\nolimits}
\def\Ext{\mathop{\mathrm{Ext}}\nolimits}
\def\SheafHom{\mathop{\mathcal{H}\!\mathit{om}}\nolimits}
\def\SheafExt{\mathop{\mathcal{E}\!\mathit{xt}}\nolimits}
\def\Sch{\mathit{Sch}}
\def\Mor{\mathop{\mathrm{Mor}}\nolimits}
\def\Ob{\mathop{\mathrm{Ob}}\nolimits}
\def\Sh{\mathop{\mathit{Sh}}\nolimits}
\def\NL{\mathop{N\!L}\nolimits}
\def\CH{\mathop{\mathrm{CH}}\nolimits}
\def\proetale{{pro\text{-}\acute{e}tale}}
\def\etale{{\acute{e}tale}}
\def\QCoh{\mathit{QCoh}}
\def\Ker{\mathop{\mathrm{Ker}}}
\def\Im{\mathop{\mathrm{Im}}}
\def\Coker{\mathop{\mathrm{Coker}}}
\def\Coim{\mathop{\mathrm{Coim}}}

% Boxtimes
%
\DeclareMathSymbol{\boxtimes}{\mathbin}{AMSa}{"02}

%
% Macros for moduli stacks/spaces
%
\def\QCohstack{\mathcal{QC}\!\mathit{oh}}
\def\Cohstack{\mathcal{C}\!\mathit{oh}}
\def\Spacesstack{\mathcal{S}\!\mathit{paces}}
\def\Quotfunctor{\mathrm{Quot}}
\def\Hilbfunctor{\mathrm{Hilb}}
\def\Curvesstack{\mathcal{C}\!\mathit{urves}}
\def\Polarizedstack{\mathcal{P}\!\mathit{olarized}}
\def\Complexesstack{\mathcal{C}\!\mathit{omplexes}}
% \Pic is the operator that assigns to X its picard group, usage \Pic(X)
% \Picardstack_{X/B} denotes the Picard stack of X over B
% \Picardfunctor_{X/B} denotes the Picard functor of X over B
\def\Pic{\mathop{\mathrm{Pic}}\nolimits}
\def\Picardstack{\mathcal{P}\!\mathit{ic}}
\def\Picardfunctor{\mathrm{Pic}}
\def\Deformationcategory{\mathcal{D}\!\mathit{ef}}
\usepackage{cleveref}
\crefname{originallemma}{original lemma}{original lemmas}
\Crefname{originallemma}{Original lemma}{Original lemmas}
\makeatletter
\newlabel{homology-lemma-first-quadrant-ss}{{0132}{source}{}{sourcepointer-1}{}}
\newlabel{homology-lemma-first-quadrant-ss@cref}{{[originallemma][0][]0132}{[1][0][]source}{}{}{}}
\makeatother
\title{A bounded above complex of flat modules on a ringed site is K-flat}
\author{The Stacks Project Authors}
\date{stacks-project:sites-cohomology.tex:local-snapshot-20261011; GNU FDL 1.2 or later}
\begin{document}
\maketitle
\noindent\textit{Editorial note.} All delivered mathematical text is the author's own, copied from the pinned Stacks Project chapter file sites-cohomology.tex; no statement or proof step was rewritten, reordered or removed.\par\medskip
\noindent\textit{Editorial note.} The proof is the author's complete argument, including the colimit reduction and the vanishing of the E\_1 terms; the author's closing phrase is retained verbatim.\par\medskip
\noindent\textit{Editorial note.} Cross-chapter restatement: the same label occurs in chapter cohomology for a ringed space. That statement is the special case of this ringed-site statement, so per the lane-s4 rule only this more general occurrence is counted here.\par\medskip
\noindent\textit{Editorial note.} The single labelled reference occurring inside the counted spans points to another chapter of the Project; it is carried in assembly.external\_references as a non-load-bearing navigation pointer recorded at the invoking line.\par\medskip
\noindent\textit{Editorial note.} The preamble is the chapter's own retained minimal preamble (sources/stacks-sites-cohomology/retained-preamble.tex), a verbatim excerpt of the native preamble: the Xy-pic setup and the native theorem declarations with the mathematical macros.\par\medskip
\noindent\textit{Allowed standard prerequisites (editorial).}\begin{itemize}
\item K-flat complex of O-modules. Same chapter, the definition of K-flatness: the total tensor product with every acyclic complex is acyclic. Supplies the property proved for a bounded above complex of flat modules.
\item Flat module on a ringed site and exactness of tensorisation. Same chapter, the definition of flatness. Supplies that each K\textasciicircum{}q is flat, which makes the E\_1 terms of the spectral sequence vanish.
\item Truncations tau\_\{<= m\} and their behaviour under tensor products. Same chapter, the truncation functors and the termwise identity L\textasciicircum{}bullet = colim\_m tau\_\{<= m\}L\textasciicircum{}bullet. Supplies the reduction from an arbitrary acyclic complex to a bounded above one, using that filtered colimits are exact.
\end{itemize}\medskip
\ifdefined\hypertarget\hypertarget{sourcepointer-1}{}\fi
\noindent\textit{Source reference.} Original lemma 0132: sources/stacks-sites-cohomology/pinned-original-source.tex, source lines 3142--3142. Cross-chapter navigation pointer invoked at proof lines 3141-3142, where the author invokes the first quadrant spectral sequence of a double complex. The target chapter is present in this repository only as another lane's snapshot, so the pointer records the invoking line; it only resolves the reference for the reader, and no bibliography citation is involved.\par\medskip
\begin{lemma}
\label{lemma-bounded-flat-K-flat}
Let $(\mathcal{C}, \mathcal{O})$ be a ringed site. A bounded above complex
of flat $\mathcal{O}$-modules is K-flat.
\end{lemma}

\begin{proof}
Let $\mathcal{K}^\bullet$ be a bounded above complex of flat
$\mathcal{O}$-modules. Let $\mathcal{L}^\bullet$ be an acyclic complex
of $\mathcal{O}$-modules. Note that
$\mathcal{L}^\bullet = \colim_m \tau_{\leq m}\mathcal{L}^\bullet$
where we take termwise colimits. Hence also
$$
\text{Tot}(\mathcal{K}^\bullet \otimes_\mathcal{O} \mathcal{L}^\bullet)
=
\colim_m \text{Tot}(
\mathcal{K}^\bullet \otimes_\mathcal{O} \tau_{\leq m}\mathcal{L}^\bullet)
$$
termwise. Hence to prove the complex on the left is acyclic it suffices
to show each of the complexes on the right is acyclic. Since
$\tau_{\leq m}\mathcal{L}^\bullet$ is acyclic this reduces us to the
case where $\mathcal{L}^\bullet$ is bounded above.
In this case the spectral sequence of
Homology, Lemma \ref{homology-lemma-first-quadrant-ss}
has
$$
{}'E_1^{p, q} = H^p(\mathcal{L}^\bullet \otimes_R \mathcal{K}^q)
$$
which is zero as $\mathcal{K}^q$ is flat and $\mathcal{L}^\bullet$ acyclic.
Hence we win.
\end{proof}

\end{document}
